% Vista editorial exacta de corpus/source/masas/04_persistencia.tex, líneas 365--590. Las líneas 1--364 ya residen exactamente en 73_rev2_electron_lectores_torres.tex.
\subsection{Persistance de l'état enrichi sous l'extension survivante et conservation de la mémoire de parcours}
Soit \(x\) un état APP, \(\tau\) son mot tritique et \(B_K\) le bloc de mémoire au niveau \(K\). L'état pertinent pour la masse prend la forme

\[
X_K=(x_K,\tau_K,\varepsilon_K,h_K,v_K,c_K,B_K,g_K),
\]
où \(\varepsilon_K\) est l'orientation du feuillet, \((h_K,v_K)\) l'orientation cellulaire, \(c_K\) la retenue et \(g_K\) la phase nonadique. Une extension est survivante lorsqu'elle forme une section compatible des cylindres successifs. Si \(C_{K,j}\), \(0\le j<729\), sont les descendants du cylindre au niveau suivant, la règle d'élagage satisfait

\[
\sum_{j=0}^{728}G_9(C_{K,j})=1.
\]
L'indice unique \(j_K\) produit la transition

\[
P_{K+1}=729P_K+j_K.
\]
Après neuf niveaux, la phase se conserve,

\[
g(K+9)=g(K),
\]
mais le bloc, le préfixe et la mémoire changent. Le retour de phase n'est donc pas un retour d'état. Une particule stable se définit par la compatibilité de la section à travers ces retours, et non par la permanence immobile d'une coordonnée visible.

Le parcours observable \(\gamma_{\rm obs}\) est l'ombre chronologique de l'extension. Le point qui parcourt un diagramme est un curseur de lecture ; la particule est la classe complète de persistance. Cette distinction permet d'interpréter correctement le mouvement quantique : il n'est pas nécessaire de supposer une petite bille suivant une trajectoire classique, mais une histoire cohérente admettant plusieurs projections compatibles.

\Needspace{7\baselineskip}
\subsection{Enregistrement dodécaphasique et signature}
La roue temporelle contient douze événements orientés. Sa décomposition causale est

\[
12=1+5+6:
\]
un événement fixe l'ancrage, cinq déterminent l'incidence visible et six forment les paires opposées qui survivent au quotient orienté. Si \(s_C^{\rm tot}\) est la somme contre-angulaire entière sur les douze événements, l'enregistrement conserve

\[
s_C:=s_C^{\rm tot}\in\mathbb Z,
\qquad
s_C^{(12)}:=\frac{s_C^{\rm tot}}6.
\]
Le premier objet est la coordonnée entière de la signature ; le second, sa lecture angulaire dodécaphasique. Le diviseur n'est pas une normalisation introduite à partir d'une unité externe. C'est l'indice qui apparaît lorsque l'on reconstruit les six paires opposées. La matrice d'incidence dodécaphasique possède une forme normale de Smith de facteurs \((1,12)\) ; le caractère est bien défini sur le quotient parce que le résidu d'orientation se mesure dans cette structure.

Telle est la convention directrice de tout le traité. Les expressions historiques conservées qui désignaient par \(s_C\) la coordonnée déjà divisée se lisent comme \(s_C^{(12)}\) ; leurs valeurs numériques et leurs témoignages demeurent intacts, mais aucune formule en vigueur ne divise deux fois la même coordonnée.

La paire d'enroulements directs est

\[
W_+=6n_A+s_C^{\rm tot},
\qquad
W_-=6n_A-s_C^{\rm tot}.
\]
Puisque \(s_C=s_C^{\rm tot}\) dans la signature publiée, cette formule coïncide avec \(W_\pm=6n_A\pm s_C\). L'exposant contre-angulaire utilise une seule fois \(s_C^{(12)}=s_C/6\) ; il ne divise pas de nouveau une coordonnée déjà normalisée.

L'enregistrement complet donne la signature

\[
\sigma(\Gamma)=
(n_A,s_C,k_{\Delta_4},b_{120},b_\zeta)\in\mathbb Z^5,
\qquad
\zeta_{270}=135\Delta_4^2.
\]
Les cinq composantes sont de types distincts. \(n_A\) compte le développement angulaire direct ; \(s_C\) enregistre la somme contre-angulaire entière ; \(s_C^{(12)}=s_C/6\) est sa lecture angulaire ; \(k_{\Delta_4}\) lit le raffinement global ; \(b_{120}\) marque le canal harmonique de hauteur 120 ; \(b_\zeta\) appartient au correcteur quadratique de hauteur 270. En particulier, \(s_{270}\) et \(\zeta_{270}\) ne sont pas le même objet : le premier est une trace impaire de l'opérateur bicouche ; le second est une coordonnée quadratique de l'enregistrement.

Le mot électronique

\[
\tau_{12}=+++---+++---
\]
a une somme linéaire nulle, mais sa signature paire n'est pas nulle :

\[
(Q_3,Q_4,Q_6,\rho_0,w)=(-12,-4,12,-6,12).
\]
Il ne faut pas non plus l'identifier à la signature brute du cycle électronique \((24,0,12,0,-12)\), ni à l'origine affine \(v_e=0\). Il s'agit de quatre niveaux de lecture distincts : mot, signature paire, signature brute et origine du caractère relatif.

\Needspace{7\baselineskip}
\subsection{Angle, contre-angle et caractère}
Pour rendre la construction autonome, nous fixons ici les trois coordonnées internes qui interviennent dans toute la loi :

\[
\alpha_{\rm HMT}
=\frac{7297352569283800997285105472380663}{10^{36}},
\]
\[
\Delta_4
=\frac{\pi_{\rm HMT}-e_{\rm HMT}\log\pi_{\rm HMT}}{270}
\left(1+\frac{\pi_{\rm HMT}}{729}\right),
\]
et, pour \(a>0\) et \(\phi>0\),

\[
H_5(a,\phi)
=\frac12\sqrt{\frac{1000a}{\phi}}-a
+\frac{9a^2}{16}-\frac{5a^3}{9}
+\frac{7a^4}{48}-\frac{a^5}{54}.
\]
Dans la suite,

\[
H_5:=H_5(\alpha_{\rm HMT},\varphi_{\rm HMT}).
\]
Ces définitions n'importent aucune constante métrologique externe : \(\pi_{\rm HMT},e_{\rm HMT},\varphi_{\rm HMT}\) sont les évaluations coinductives du lecteur HMT et \(\alpha_{\rm HMT}\) est le résultat rationnel de la clôture dodécaphasique. La constante interne \(\alpha_{\rm HMT}\) précède le contre-angle. La branche régionale détermine d'abord

\[
A=1000\alpha_{\rm HMT},
\qquad
D_A=\exp\!\left(-\frac{100\pi\alpha_{\rm HMT}}9\right),
\]
\[
\mathcal C_\pi(\alpha_{\rm HMT})
=169\alpha_{\rm HMT}^{6}
+\frac{D_A\alpha_{\rm HMT}^{7}}
       {1-D_A\alpha_{\rm HMT}},
\qquad
y_*=\frac{\pi}{90}(H_5+\alpha_{\rm HMT})
-\mathcal C_\pi(\alpha_{\rm HMT}),
\]
\[
C^*=\frac{180}{\pi}y_*.
\]
Ici, \(A\) et \(C^*\) sont des coordonnées angulaires déjà construites, et non des grandeurs d'action. Si

\[
\alpha_\angle:=\alpha_{\rm HMT}\,{}^\circ,
\qquad
\eta_{\rm ret}^{(\deg)}:=\frac{C^*}{2}-\alpha_\angle,
\]
alors

\[
C^*=2\bigl(\eta_{\rm ret}^{(\deg)}+\alpha_\angle\bigr)
\]
est l'identité inverse de retour, et non une seconde règle de sélection du contre-angle. Cette écriture évite de confondre la mantisse angulaire \(\eta_{\rm ret}^{(\deg)}\) avec la grandeur dimensionnelle \(\hbar_{\rm ret}\). Les cartes en radians s'obtiennent par multiplication par \(\pi/180\). Les canaux conjugués sont

\[
q_\pm=\exp\left[-\frac\pi{180}(A\pm C^*)\right].
\]
Le caractère central d'une signature est

\[
\chi_0(\sigma)=
\exp\left(
n_AA_{\rm rad}
+\frac{s_C}{6}C^*_{\rm rad}
+k_{\Delta_4}\Delta_4
+b_\zeta\zeta_{270}
\right).
\]
La coordonnée \(b_{120}\) demeure dans la signature, mais sa contribution est incorporée dans la famille de tours. Cette convention empêche de compter deux fois la même hauteur.

Le caractère est adimensionnel. Sa fonction n'est pas de devenir une constante d'action, mais d'agir par homothéties positives sur une droite de masse déjà typée. Pour une paire d'états \(X,Y\),

\[
\frac{m_Y^{(0)}}{m_X^{(0)}}
=\chi_{\rm full}(\sigma_Y-\sigma_X,\eta_Y-\eta_X).
\]
Cette formule sépare les rapports de base des masses du choix commun de carte absolue. Le quotient bêta complet incorpore ensuite \(R_{\rm act}^{\kappa_Y-\kappa_X}\).

\Needspace{7\baselineskip}
\subsection{Semi-groupe global des tours de masse et développement harmonique convergent}
Soit \(R^2=I\), avec les projecteurs

\[
P_\pm=\frac{I\pm R}{2},
\qquad
x=\frac{\pi A}{180},
\qquad
y=\frac{\pi C^*}{180},
\]
et

\[
T=e^{-xI-yR}=q_+P_++q_-P_-.
\]
Les traces paire et impaire sont

\[
c_n=\operatorname{Tr}(T^n)=q_-^n+q_+^n
=2e^{-nx}\cosh(ny),
\]
\[
s_n=-\operatorname{Tr}(RT^n)=q_-^n-q_+^n
=2e^{-nx}\sinh(ny).
\]
Elles satisfont

\[
c_n^2-s_n^2=4e^{-2nx},
\]
\[
c_{m+n}=\frac{c_mc_n+s_ms_n}{2},
\qquad
s_{m+n}=\frac{s_mc_n+c_ms_n}{2},
\]
et la récurrence linéaire d'ordre deux

\[
X_{n+2}=c_1X_{n+1}-e^{-2x}X_n,
\qquad X\in\{c,s\},
\qquad c_0=2,\quad s_0=0.
\]
L'ensemble des hauteurs ne se réduit pas à huit corrections visibles. C'est le semi-groupe

\[
\langle90,120\rangle
=\{90,120\}\cup\{30r:r\ge6\}.
\]
Les hauteurs \(90,120,180,210,240,270,360,420\) sont les premiers témoins généalogiques utiles, et non le domaine entier. Le caractère raffiné est

\[
\log\chi_{\rm full}(\sigma,\eta)
=\log\chi_0(\sigma)
+\sum_{h\in\langle90,120\rangle}N_h(\sigma,\eta)s_h,
\]
avec

\[
N_h(\sigma,\eta)=\eta_h+b_{120}\delta_{h,120}.
\]
Ainsi, \(s_{120}\) apparaît une seule fois, et à la hauteur 270, \(N_{270}s_{270}\) et \(b_\zeta\zeta_{270}\) restent distincts. La convergence absolue exige

\[
\sum_h |N_hs_h|<\infty.
\]
La complétude du semi-groupe garantit la capacité de représentation des rapports positifs ; elle ne choisit pas, à elle seule, les coefficients physiques d'une espèce. La règle prospective de sélection, lorsqu'elle intervient, doit procéder du parcours, de l'enregistrement généalogique, de la retenue, de la mémoire et de l'enroulement avant l'introduction de la masse externe.

\Needspace{7\baselineskip}
